\section*{अध्याय \ref{ch:g8:plastics} --- प्लास्टिक और संश्लिष्ट सामग्रियाँ}

\begin{solution}{exo:g8:plastics:1}
प्राकृतिक: सूत, ऊन। कृत्रिम: कागज़, विस्कोस। संश्लिष्ट: नाइलॉन,
पॉलिस्टाइरीन।
\end{solution}

\begin{solution}{exo:g8:plastics:2}
कच्चे तेल (या प्राकृतिक गैस) से।
\end{solution}

\begin{solution}{exo:g8:plastics:3}
PET (पॉलिएथिलीन टेरेफ़्थेलेट): पानी और सोडा की बोतलें।
\end{solution}

\begin{solution}{exo:g8:plastics:4}
\omterm{def:g8:plastics:thermoplastic}{थर्मोप्लास्टिक} हर बार गरम करने पर नरम होता है और फिर से ढाला जा सकता है;
\omterm{def:g8:plastics:thermoplastic}{थर्मोसेट} पहली बार में ही हमेशा के लिए जम जाता है और नरम होने के बजाय झुलस जाता है।
\end{solution}

\begin{solution}{exo:g8:plastics:5}
पॉलिप्रोपिलीन (PP), एक \omterm{def:g8:plastics:thermoplastic}{थर्मोप्लास्टिक}: उसे पिघलाकर फिर से ढाला जा सकता है।
\end{solution}

\begin{solution}{exo:g8:plastics:6}
त्रिभुज केवल \omterm{def:g8:plastics:plastic}{प्लास्टिक} का नाम बताता है, ताकि छँटाई में मदद मिले। वस्तु का
\omterm{def:g5:raw-materials-and-recycling:recycling}{पुनर्चक्रण} इस पर निर्भर है कि वह इकट्ठी की जाए और कोई संयंत्र उस \omterm{def:g8:plastics:plastic}{प्लास्टिक}
का \omterm{def:g5:raw-materials-and-recycling:recycling}{पुनर्चक्रण} करता हो।
\end{solution}

\begin{solution}{exo:g8:plastics:7}
कड़ाही गरम होती है: \omterm{def:g8:plastics:thermoplastic}{थर्मोप्लास्टिक} नरम हो जाएगा, \omterm{def:g8:plastics:thermoplastic}{थर्मोसेट} नहीं होता।
\end{solution}

\begin{solution}{exo:g8:plastics:8}
$\frac{353}{460} \approx 0.77$: लगभग \qty{77}{\percent}।
\end{solution}

\begin{solution}{exo:g8:plastics:9}
PVC में क्लोरीन है: उसे \omterm{def:g4:what-a-fire-needs:burning}{जलाने} पर हाइड्रोजन क्लोराइड निकलता है, जो एक
संक्षारक गैस है, और साथ में बगीचे की आग के \omterm{def:g8:combustion-and-fuels:complete}{अपूर्ण दहन} से कालिख और कार्बन
मोनोऑक्साइड भी।
\end{solution}

\begin{solution}{exo:g8:plastics:10}
$300 \times 36 = 10\,800$ बोतलें; $10\,800 \times 25 = 270\,000$\,g, यानी
\qty{270}{kg}।
\end{solution}

\begin{solution}{exo:g8:plastics:11}
कार्बन डाइऑक्साइड और पानी। उसका कार्बन कच्चे तेल से आता है, जो लाखों वर्षों
से ज़मीन के नीचे जमा था: उसे \omterm{def:g4:what-a-fire-needs:burning}{जलाने} से \omterm{def:g4:air-a-mixture-of-gases:air}{वायु} में कार्बन डाइऑक्साइड जुड़ती है,
जैसे \omterm{def:g8:combustion-and-fuels:fossil}{जीवाश्म ईंधन} \omterm{def:g4:what-a-fire-needs:burning}{जलाने} से।
\end{solution}

\begin{solution}{exo:g8:plastics:12}
बोतल सड़ती नहीं: धूप और लहरें उसे छोटे से छोटे टुकड़ों में तोड़ती जाती हैं।
नन्हे टुकड़े पानी में तैरते या डूबते हैं, और मछलियाँ उन्हें अपने भोजन के साथ
निगल लेती हैं।
\end{solution}

\begin{solution}{pb:g8:plastics:1}
\textbf{1.} ढाँचे: 1 (PET)। ढक्कन: 2 (HDPE) या 5 (PP)।

\textbf{2.} संश्लिष्ट: वे रसायनज्ञों द्वारा बनाए गए पदार्थों से, ज़्यादातर
तेल से, बनते हैं।

\textbf{3.} हाँ: PET एक \omterm{def:g8:plastics:thermoplastic}{थर्मोप्लास्टिक} है।

\textbf{4.} हर \omterm{def:g8:plastics:plastic}{प्लास्टिक} का \omterm{def:g5:raw-materials-and-recycling:recycling}{पुनर्चक्रण} अलग से होता है: साथ पिघलाने पर \omterm{def:g8:plastics:plastic}{प्लास्टिकों}
का \omterm{def:g6:pure-substances-and-mixtures:mixture}{मिश्रण} घटिया सामग्री देता है।

\textbf{5.} $0.80 \times 1000 = \qty{800}{kg}$।

\textbf{6.} ढक्कन: $0.12 \times 1000 = \qty{120}{kg}$। लेबल: $0.05
\times 1000 = \qty{50}{kg}$।

\textbf{7.} $0.03 \times 1000 = \qty{30}{kg}$ गंदगी और द्रव।

\textbf{8.} $80 + 12 + 5 = \qty{97}{\percent}$।

\textbf{9.} $0.025 \times 800 = \qty{20}{kg}$।

\textbf{10.} टुकड़े कच्चे तेल को और नया PET बनाने की ऊर्जा को बचाते हैं,
और कचरे को एक उपयोगी वस्तु में बदल देते हैं।

\textbf{11.} कार्बन डाइऑक्साइड (और, अगर \omterm{def:g4:what-a-fire-needs:burning}{दहन} ठीक से न हो, तो कार्बन
मोनोऑक्साइड और कालिख)।

\textbf{12.} प्रति गट्ठर $800 - 20 = \qty{780}{kg}$ साफ़ PET टुकड़े।
\end{solution}
